Titel:

   Metastability: a universal approach through potential theory


Abstract: In these lectures I will give an introduction and overview of the
  achievement and challanges of the so-called potential theoretic approach to 
metastablility. There will be three parts:

Part 1: Potential theory for (reversible) 
Markov processes and the abstract definition of metastables sets in terms of 
capacities. Emergence of universal features.
Part 2: Variational principles and the computation of capacities. 
Blocking methods and the production of entropy. Examples.
Part 3: Challanges in high dimesnions, coupling methods and regularity. 

The talk by Florent Barret will show a nice application of the general theory 
in the setting of stochastic spdes. 

Reference: A. Bovier, Metastability,  in Methods of Contemporary Statistical 
Mechanics, R. Kotecky, Ed., LNM 1970, Springer, 2009.